\documentclass[10pt,a4paper]{amsart}
\usepackage[margin=19mm]{geometry}
\usepackage{amsmath,amssymb,mathtools}
\usepackage[scr=rsfs,cal=euler]{mathalfa}
\usepackage{xfrac}
\usepackage[unicode,hidelinks,bookmarks=false]{hyperref}
\newcommand{\ie}{i.e.\ }
\newcommand{\cf}{cf.\ }
\newcommand{\resp}{respectively\ }
\newcommand{\Subsection}{\subsection}
\providecommand{\nobreakdash}{\nobreak\hbox{-}}
\setlength{\emergencystretch}{2em}
\multlinegap0pt
\usepackage{amssymb}
\usepackage{enumerate}
\newtheorem{theorem}{Theorem}
\newtheorem{lemma}{Lemma}
\usepackage{cleveref}
\crefname{theorem}{Theorem}{Theorems}
\crefname{lemma}{Lemma}{Lemmas}
\providecommand{\cM}{\mathcal{M}}
\title{Uniqueness for 2-Intersecting Families of \\ Permutations and Perfect Matchings}
\author{Gilad Chase, Neta Dafni, Yuval Filmus, Nathan Lindzey}
\date{Source: arXiv:2210.00245v2 (CC BY 4.0)}
\begin{document}
\maketitle

\noindent\emph{Source and licence.} Uniqueness for 2-Intersecting Families of \\ Permutations and Perfect Matchings. Version arXiv:2210.00245v2 (CC BY 4.0), distributed by arXiv under the Creative Commons Attribution 4.0 International licence, http://creativecommons.org/licenses/by/4.0/, as recorded in the arXiv licence metadata for this exact version and shown on the abstract page https://arxiv.org/abs/2210.00245v2. Full licence notice: \texttt{licenses/arxiv-2210.00245-CC-BY-4.0.txt}.\par\smallskip

\noindent\textit{Editorial scope.} Counted unit: the article's theorem thm:main-pms (source0.tex lines 106--108). The article's complete original proof of it is delivered with it (source0.tex lines 763--807, one \texttt{proof} environment of 45 lines, which the article places away from the statement). No mathematics is written, repaired or re-derived: the statement and the proof are the article's own lines, in the article's own notation, and no retained line is substituted or re-flowed.  Retained uncounted as the source-local context the proof needs: lines 91--96; lines 436--438; lines 490--492; lines 742--750; lines 861--861.  Nothing was omitted from any retained span, so the package carries no omission record.  No retained line contains a citation, so the wrapper carries no bibliography and no bibliography entry.  No retained line contains a figure environment, an image-inclusion command or drawing code, so no image is delivered and the package declares no asset dependency.  Editorial formatting only, in addition: an AMS \texttt{amsart} preamble that restates the article's own declarations and macros used by the retained lines, namely the environments \texttt{theorem}, \texttt{lemma} on separate counters, together with the standard packages those lines need.  The retained environments print in this document as Theorem 1, Lemma 1 in order of appearance, whereas the article prints the same items with the numbering of its own full document.  The complete native source of the article as distributed in the arXiv e-print for this exact version is bundled unchanged as \texttt{sources/131467/source0.tex}.
\begin{theorem} \label{thm:main-sym}
Let $n \geq 2$. If $\mathcal{F}$ is a $2$-intersecting subset of $S_n$ of size $(n-2)!$, then there exist $i_1 \neq i_2$ and $j_1 \neq j_2$ such that
\[
 \mathcal{F} = \{ \alpha \in S_n : \alpha(i_1) = j_1 \text{ and } \alpha(i_2) = j_2 \}.
\]
\end{theorem}

\begin{theorem} \label{thm:FMS}
If $n \geq 3$ then any $2$-intersecting family $\mathcal{F} \subseteq \cM_{2n}$ contains at most $(2n-5)!!$ perfect matchings. Furthermore, if $f\colon \cM_{2n} \to \{0,1\}$ is the characteristic vector of a $2$-intersecting family of size $(2n-5)!!$, then $\deg f \leq 2$.
\end{theorem}

\begin{lemma} \label{lem:C-bound-pms}
If $n \geq 8$ and $f\colon \cM_{2n} \to \{0,1\}$ has degree at most~$2$, then $C(f) \leq n-2$.
\end{lemma}

\begin{lemma} \label{lem:main-ub-pms}
Let $\mathcal{F} \subseteq \cM_{2n}$ be a maximal $2$-intersecting family whose characteristic function $f$ has degree at most~$2$, and suppose that $C(f) \leq n-2$. Let $C_1(f)$ the maximum certificate complexity of any $m \in \mathcal{F}$.

The family $\mathcal{F}$ is contained in the union of at most $T$ many $C_1(f)$-cosets, where
\[
 T = \frac{2\lfloor C_1(f)/2 \rfloor (C_1(f) - 1)!}{2^{\lfloor C_1(f)/2 \rfloor}}.
\]
Moreover, $C_1(f) \geq 2$.
\end{lemma}

\subsection{Code for Theorem \ref{thm:main-pms}} \label{apx:main-pms}

\begin{theorem} \label{thm:main-pms}
Let $n \geq 2$. If $\mathcal{F}$ is a $2$-intersecting subset of $\cM_{2n}$ of size $(2n-5)!!$, then there exist distinct $i_1,j_1,i_2,j_2$ such that $\mathcal{F}$ consists of all perfect matchings containing the edges $\{i_1,j_1\},\{i_2,j_2\}$.
\end{theorem}

\begin{proof}[Proof of \Cref{thm:main-pms}]
Let $\mathcal{F}$ be a $2$-intersecting subset of $\cM_{2n}$ of size $(2n-5)!!$, where $n \geq 2$, and let $f$ be its characteristic function. Our goal is to show that
\[
 \mathcal{F} = \{ m \in \cM_{2n} : m(i_1) = j_1 \text{ and } m(i_2) = j_2 \}
\]
for some distinct $i_1,j_1,i_2,j_2 \in [2n]$.

Suppose first that $n \geq 8$. Since $n \geq 3$, \Cref{thm:FMS} shows that $\mathcal{F}$ is maximal and that $\deg f \leq 2$. \Cref{lem:C-bound-pms} shows that $C_1(f) \leq n-2$. \Cref{lem:main-ub-pms} implies that $C_1(f) \geq 2$ and  $|\mathcal{F}| \leq T (2n - 2C_1(f) - 1)!!$, where $T = \lfloor C_1(f)/2 \rfloor (C_1(f) - 1)!/2^{\lfloor C_1(f)/2 \rfloor}$. Therefore
\[
 \frac{(2n-5)!!}{(2n - 2C_1(f) - 1)!!} \leq \frac{2\lfloor C_1(f)/2 \rfloor (C_1(f) - 1)!}{2^{\lfloor C_1(f)/2 \rfloor}}.
\]
The left-hand side is clearly increasing in $n$. Since $n \geq C_1(f) + 2$, it follows that
\[
 \frac{(2C_1(f))!}{3 \cdot 2^{C_1(f)} C_1(f)!} =
 \frac{(2C_1(f) - 1)!!}{3} \leq \frac{2\lfloor C_1(f)/2 \rfloor (C_1(f) - 1)!}{2^{\lfloor C_1(f)/2 \rfloor}},
\]
and so 
\[
 (2C_1(f))! \leq 3 \cdot 2^{\lceil C_1(f)/2 \rceil} C_1(f)! (2\lfloor C_1(f)/2 \rfloor) (C_1(f) - 1)!.
\]

If $C_1(f)$ is even then this reads
\[
 (2C_1(f))! \leq 3 \cdot 2^{C_1(f)/2} C_1(f)!^2.
\]
When $C = C_1(f)$ increases by $1$, the left-hand side increases by a factor of $(2C+2)(2C+1)$, while the right-hand side increases by a factor of $\sqrt{2} (C+1)^2$, which is smaller for $C \geq 2$. When $C = 3$, the inequality fails. Therefore in this case, $C_1(f) = 2$.

If $C_1(f)$ is odd then the inequality reads
\[
 (2C_1(f))! \leq 3 \cdot 2^{(C_1(f)+1)/2} C_1(f)! (C_1(f) - 1)(C_1(f) - 1)! < 3 \cdot 2^{(C_1(f)+1)/2} C_1(f)!^2.
\]
As before, if we increase $C = C_1(f)$ by $1$, the left-hand side increases by a larger factor than the (far) right-hand side. When $C = 3$, the inequality fails, and so this case cannot happen.

Now let $m \in \mathcal{F}$ be arbitrary. Since $C_1(f) = 2$, there is a certificate of the form $\bigl\{\{i_1,j_1\},\{i_2,j_2\}\bigr\}$ for $m$. The number of permutations satisfying this certificate is $(2n-5)!!$, and we conclude that $\mathcal{F}$ consists of all perfect matchings satisfying the certificate, completing the proof in this case.

\smallskip

In order to complete the proof, we consider the case $n \leq 7$. It suffices to show that for $n \in \{2,3,4,5,6,7\}$, any $2$-intersecting family of size $(2n-5)!!$ which contains the perfect matching $\{1,n+1\},\ldots,\{n+2n\}$ is of the form
\[
 \{ m \in \cM_{2n} : m(i_1) = n + i_1, m(i_2) = n + i_2 \}
\]
for some $i_1 \neq i_2 \in [n]$.

If $n = 2$ or $n = 3$ then $(2n-5)!! = 1$, and the claim can be verified directly. In order to verify the remaining cases $n \in \{4,5,6,7\}$, we use exhaustive search just as in the proof of \Cref{thm:main-sym}. The relevant code appears in \Cref{apx:main-pms}.
\end{proof}

\end{document}
